\chapter{Lines and Polygons}\label{ch:g4:geometry}

With a ruler, a set square and a compass, geometry becomes a craft.
This chapter teaches the two special positions of lines ---
\omterm{def:g4:geometry:def}{perpendicular} and \omterm{def:g4:geometry:def}{parallel} --- at the tool level, and takes a tour of
the \omterm{def:g4:geometry:polygon}{polygon} family. (The theory of these notions comes in
\cref{ch:g6:lines}.)

\section{Perpendicular and parallel}

\begin{definition}[Perpendicular, parallel]\label{def:g4:geometry:def}
Two lines are \emph{perpendicular}\index{perpendicular lines} when
they cross at a \omterm{def:g3:shapes:rightangle}{right angle} (checked with the set square, marked with
a little square). Two lines are \emph{parallel}\index{parallel lines}
when they run in the same direction and never cross --- like the two
rails of a track, always at the same distance from each other.
\end{definition}

\begin{method}[Drawing with the set square]\label{met:g4:geometry:draw}
\emph{\omterm{def:g4:geometry:def}{Perpendicular} to a line $d$ through a point $P$}: lay one edge
of the \omterm{def:g3:shapes:rightangle}{right angle} along $d$, slide until the other edge reaches $P$,
draw.

\emph{\omterm{def:g4:geometry:def}{Parallel} to $d$ through $P$}: lay the set square's edge along
$d$, press the ruler against another edge as a rail, slide the set
square along the ruler up to $P$, draw.
\end{method}

\begin{omfigure}
\begin{tikzpicture}[scale=0.85]
  \draw[thick] (-0.4,0) -- (4.2,0) node[below, font=\small] {$d$};
  \draw[omDef, thick] (1.8,-0.9) -- (1.8,2.3);
  \draw[thin] (1.8,0) ++(-0.25,0) -- ++(0,0.25) -- ++(0.25,0);
  \fill (1.8,1.6) circle (1.6pt) node[right, font=\small] {$P$};
  \node[below, font=\small] at (1.9,-1.1) {perpendicular through $P$};
  \begin{scope}[xshift=7.4cm]
  \draw[thick] (-0.4,0) -- (4.2,0) node[below, font=\small] {$d$};
  \draw[omThm, thick] (-0.4,1.6) -- (4.2,1.6);
  \fill (1.6,1.6) circle (1.6pt) node[above, font=\small] {$P$};
  \draw[Stealth-Stealth, thin] (3.4,0) -- (3.4,1.6)
    node[midway, right, font=\small] {same gap everywhere};
  \node[below, font=\small] at (1.9,-1.1) {parallel through $P$};
  \end{scope}
\end{tikzpicture}

{\small The two constructions. \omterm{def:g4:geometry:def}{Parallel lines} keep a constant gap:
measuring the gap at two places is a good check.}
\end{omfigure}

\begin{example}[Spotting them in figures]\label{ex:g4:geometry:spot}
In a rectangle, each side is \omterm{def:g4:geometry:def}{perpendicular} to its two neighbors and
\omterm{def:g4:geometry:def}{parallel} to the opposite side. Grid paper is a ready-made network of
\omterm{def:g4:geometry:def}{parallels} and \omterm{def:g4:geometry:def}{perpendiculars} --- that is what makes it so
comfortable to draw on.
\end{example}

\section{Polygons}

\begin{definition}[Polygon]\label{def:g4:geometry:polygon}
A \emph{polygon}\index{polygon} is a closed figure whose border is
made of straight segments, its \emph{sides}; the corner points are
the \emph{vertices}\index{vertex}. Polygons are named by their number of sides:
\emph{triangle} ($3$), \emph{quadrilateral}\index{quadrilateral} ($4$),
\emph{pentagon}\index{pentagon} ($5$), \emph{hexagon}\index{hexagon} ($6$),
\emph{octagon}\index{octagon}
($8$).
\end{definition}

\begin{omfigure}
\begin{tikzpicture}[scale=0.62]
  \draw[thick, fill=omDef!10] (0,0) -- (2,0) -- (1.2,1.7) -- cycle;
  \node[below, font=\small] at (1,-0.3) {triangle};
  \begin{scope}[xshift=3.4cm]
  \draw[thick, fill=omThm!10] (0,0) -- (2.2,0) -- (2.5,1.5) --
    (0.5,1.8) -- cycle;
  \node[below, font=\small] at (1.25,-0.3) {quadrilateral};
  \end{scope}
  \begin{scope}[xshift=7.4cm]
  \draw[thick, fill=omMeth!10] (90:1.1) -- (162:1.1) -- (234:1.1) --
    (306:1.1) -- (18:1.1) -- cycle;
  \node[below, font=\small] at (0,-1.4) {pentagon};
  \end{scope}
  \begin{scope}[xshift=10.9cm]
  \draw[thick, fill=omProp!10] (0:1.1) -- (60:1.1) -- (120:1.1) --
    (180:1.1) -- (240:1.1) -- (300:1.1) -- cycle;
  \node[below, font=\small] at (0,-1.4) {hexagon};
  \end{scope}
\end{tikzpicture}

{\small The \omterm{def:g4:geometry:polygon}{polygon} family (the \omterm{def:g4:geometry:polygon}{pentagon} and \omterm{def:g4:geometry:polygon}{hexagon} drawn here are
\emph{regular}: equal sides and equal corners). A circle is not a
\omterm{def:g4:geometry:polygon}{polygon}: its border is not made of segments.}
\end{omfigure}

\begin{example}[Describing a polygon]\label{ex:g4:geometry:describe}
A useful game: describe a figure so precisely that a classmate can
redraw it without seeing it. ``A \omterm{def:g4:geometry:polygon}{quadrilateral} $ABCD$ with
$AB = 5$ cm, the angle at $A$ right, $AD = 3$ cm\,\dots'' ---
vocabulary is what makes the message work.
\end{example}

\section{The circle and the compass}

\begin{definition}[Circle]\label{def:g4:geometry:circle}
The \emph{circle} of center $O$ and \omterm{def:g3:shapes:shapes}{radius} $3$ cm is the set of all
points at exactly $3$ cm from $O$. The compass draws it: spike on
$O$, opening $3$ cm, turn. A \emph{diameter}\index{diameter} is a segment through
the center joining two points of the circle; it measures twice the
\omterm{def:g3:shapes:shapes}{radius}.
\end{definition}

\begin{example}[The compass as a distance-keeper]\label{ex:g4:geometry:compass}
The compass does more than draw circles: it \emph{transports
distances}. To mark all points at $4$ cm from a point $A$, draw the
circle of center $A$, \omterm{def:g3:shapes:shapes}{radius} $4$ cm; to compare two segments without
a ruler, open the compass on one and test it on the other.
\end{example}

\begin{method}[Reproducing a figure on a grid]\label{met:g4:geometry:reproduce}
\begin{enumerate}
  \item Choose a starting \omterm{def:g4:geometry:polygon}{vertex} and its position on the blank grid;
  \item travel along the figure \omterm{def:g4:geometry:polygon}{vertex} by \omterm{def:g4:geometry:polygon}{vertex}, counting the grid
        steps (``$3$ right, $2$ up'');
  \item join the \omterm{def:g4:geometry:polygon}{vertices} in order, and compare with the model:
        same lengths, same directions.
\end{enumerate}
\end{method}

\section{Exercises}

\begin{exercise}[$\star$]\label{exo:g4:geometry:1}
Draw a line $d$ and a point $P$ not on it. Construct the
\omterm{def:g4:geometry:def}{perpendicular} to $d$ through $P$, then the \omterm{def:g4:geometry:def}{parallel} to $d$ through
$P$ (with the ruler-rail method). Mark the \omterm{def:g3:shapes:rightangle}{right angle}.
\end{exercise}

\begin{exercise}[$\star$]\label{exo:g4:geometry:2}
Look at the capital letters E, H, N, T, Z. In each, find a pair of
\omterm{def:g4:geometry:def}{parallel} strokes and a pair of \omterm{def:g4:geometry:def}{perpendicular} strokes when they
exist.
\end{exercise}

\begin{exercise}[$\star$]\label{exo:g4:geometry:3}
Name each \omterm{def:g4:geometry:polygon}{polygon}: $6$ sides; $4$ sides; $3$ sides; $8$ sides;
$5$ sides. How many \emph{\omterm{def:g4:geometry:polygon}{vertices}} does each one have?
\end{exercise}

\begin{exercise}[$\star$]\label{exo:g4:geometry:4}
Draw any \omterm{def:g4:geometry:polygon}{pentagon} (not necessarily regular) and label its \omterm{def:g4:geometry:polygon}{vertices}
$A$, $B$, $C$, $D$, $E$. How many sides does it have? Draw its
diagonals from $A$ (segments joining $A$ to non-neighbor \omterm{def:g4:geometry:polygon}{vertices}):
how many are there?
\end{exercise}

\begin{exercise}[$\star$]\label{exo:g4:geometry:5}
Draw a circle of center $O$ and \omterm{def:g3:shapes:shapes}{radius} $4$ cm; a \omterm{def:g4:geometry:circle}{diameter} $[AB]$;
and a point $M$ of the circle not on $[AB]$. Give the lengths $OA$,
$OM$ and $AB$ without measuring.
\end{exercise}

\begin{exercise}[$\star$]\label{exo:g4:geometry:6}
With the compass as a distance-keeper
(\cref{ex:g4:geometry:compass}), compare the sides of a triangle
drawn by a classmate and decide whether it has two equal sides.
\end{exercise}

\begin{exercise}[$\star$]\label{exo:g4:geometry:7}
On grid paper, reproduce this path and close it into a \omterm{def:g4:geometry:polygon}{polygon}:
start at a point $A$; go $4$ right; $2$ up; $3$ left; $2$ up;
then back to $A$ in a straight line. What is the name of the \omterm{def:g4:geometry:polygon}{polygon}
you get (count its sides)?
\end{exercise}

\begin{exercise}[$\star$]\label{exo:g4:geometry:8}
True or false? Explain with a drawing.
``Two \omterm{def:g4:geometry:def}{perpendicular lines} always cross.''
``Two \omterm{def:g4:geometry:def}{parallel lines} never cross.''
``A line can be \omterm{def:g4:geometry:def}{parallel} to one line and \omterm{def:g4:geometry:def}{perpendicular} to
another.''
\end{exercise}

\begin{exercise}[$\star\star$]\label{exo:g4:geometry:9}
Draw a line $d$, then a line $e$ \omterm{def:g4:geometry:def}{perpendicular} to $d$, then a line
$f$ \omterm{def:g4:geometry:def}{perpendicular} to $e$. Look at $d$ and $f$: what do you notice?
(This observation becomes a theorem in \cref{ch:g6:lines}.)
\end{exercise}

\begin{exercise}[$\star\star$]\label{exo:g4:geometry:10}
Describe the figure below in words for a classmate, without showing
it: a square of side $4$ cm, with a circle of \omterm{def:g3:shapes:shapes}{radius} $2$ cm centered
at the middle of the square (touching the four sides). Write the
description as numbered instructions.
\end{exercise}

\begin{exercise}[$\star\star$]\label{exo:g4:geometry:11}
How many diagonals does a \omterm{def:g4:geometry:polygon}{quadrilateral} have? A \omterm{def:g4:geometry:polygon}{pentagon}? A \omterm{def:g4:geometry:polygon}{hexagon}?
Draw the three cases, count carefully, and describe the pattern you
see.
\end{exercise}
